\documentclass[12pt,a4paper]{article}

% Packages
\usepackage[utf8]{inputenc}
\usepackage[german]{babel}
\usepackage{geometry}
\usepackage{hyperref}
\usepackage{longtable}
\usepackage{parskip}

\geometry{a4paper, margin=1in}

% Title Page
\title{Lastenheft \\ \large Projektname}
\author{Auftraggeber / Dein Name}
\date{\today}

\begin{document}

\maketitle
\tableofcontents
\newpage

%--------------------------------------------------
\section{Einleitung}
\subsection{Ziel des Dokuments}
Dieses Lastenheft beschreibt die fachlichen Anforderungen an ein System zur automatisierten Lieferterminplanung und -bestätigung. Ziel ist es, den bestehenden manuellen Prozess zu digitalisieren und effizienter zu gestalten.

\subsection{Projektübersicht}
Im aktuellen Prozess werden Liefertermine telefonisch durch Disponenten mit Kunden abgestimmt. Dieser Ansatz ist zeitaufwendig, fehleranfällig und nicht skalierbar.

Ziel des Projekts ist die Entwicklung eines prototypischen Systems, das diesen Prozess automatisiert. Das System soll Liefertermine planen, dem Kunden vorschlagen und dessen Rückmeldung automatisch verarbeiten.
%--------------------------------------------------
\section{Ist-Zustand}
Der aktuelle Prozess basiert auf manueller Kommunikation per Telefon zwischen Disponenten und Kunden.

Probleme:
\begin{itemize}
    \item Hoher Zeitaufwand für Terminabstimmungen
    \item Fehleranfälligkeit durch manuelle Eingaben
    \item Keine transparente Nachverfolgung von Statusänderungen
    \item Schlechte Skalierbarkeit bei wachsender Auftragszahl
\end{itemize}

%--------------------------------------------------
\section{Zielbestimmung}
\subsection{Muss-Ziele}
\begin{itemize}
    \item Automatische Planung von Lieferterminen
    \item Digitale Benachrichtigung von Kunden (E-Mail, WhatsApp oder Weblink)
    \item Verarbeitung von Kundenreaktionen (Bestätigung, Ablehnung, keine Antwort)
    \item Automatische Anpassung der Planung bei Ablehnung oder fehlender Rückmeldung
    \item Speicherung aller Statusänderungen im System
\end{itemize}

\subsection{Kann-Ziele}
\begin{itemize}
    \item Integration zusätzlicher Kommunikationskanäle
    \item Erweiterte Monitoring- und Reporting-Funktionen
    \item Optimierung der Tourenplanung durch intelligente Algorithmen
\end{itemize}

\subsection{Abgrenzungskriterien}
\begin{itemize}
    \item Keine vollständige Integration in bestehende ERP-Systeme
    \item Keine Entwicklung eines produktiven Systems (nur Prototyp)
    \item Keine komplexe Routenoptimierung
\end{itemize}

%--------------------------------------------------
\section{Produkt-Einsatz}
\subsection{Anwendungsbereiche}
Das System wird im Bereich der Disposition zur Planung und Bestätigung von Lieferterminen eingesetzt.

\subsection{Zielgruppen}
\begin{itemize}
    \item Disponenten
    \item Kunden
\end{itemize}

\subsection{Betriebsbedingungen}
\begin{itemize}
    \item Webbasierte Anwendung
    \item Zugriff über Desktop und mobile Endgeräte
    \item Online-Betrieb erforderlich
\end{itemize}

%--------------------------------------------------
\section{Funktionale Anforderungen}
\begin{longtable}{|l|p{10cm}|}
\hline
ID & Beschreibung \\
\hline
F1 & System plant automatisch einen Liefertermin basierend auf verfügbaren Zeitfenstern \\
F2 & System sendet Terminvorschläge an Kunden (E-Mail, WhatsApp oder Weblink) \\
F3 & Kunde kann Termin bestätigen, ablehnen oder nicht reagieren \\
F4 & System verarbeitet Kundenreaktionen automatisch \\
F5 & Bei Bestätigung wird der Auftrag fix in die Tour eingeplant \\
F6 & Bei Ablehnung werden alternative Zeitfenster geprüft \\
F7 & Falls keine Alternative vorhanden ist, wird eine neue Tour geprüft \\
F8 & System speichert alle Statusänderungen \\
F9 & Benutzer können Prozesse starten und überwachen (Frontend) \\
\hline
\end{longtable}

%--------------------------------------------------
\section{Nicht-funktionale Anforderungen}
\begin{itemize}
    \item Performance: Reaktionszeit unter 2 Sekunden für Benutzeraktionen
    \item Sicherheit: DSGVO-konforme Verarbeitung personenbezogener Daten
    \item Usability: Intuitive Benutzeroberfläche
    \item Skalierbarkeit: Erweiterbar für größere Datenmengen
\end{itemize}

%--------------------------------------------------
\section{Daten}
\subsection{Zu speichernde Daten}
\begin{itemize}
    \item Auftragsdaten
    \item Kundendaten
    \item Liefertermine und Zeitfenster
    \item Statusinformationen
\end{itemize}

\subsection{Datenschutz}
Alle personenbezogenen Daten müssen gemäß DSGVO verarbeitet und gespeichert werden.

%--------------------------------------------------
\section{Systemgrenzen}
Das System interagiert mit einem DISPO-System, übernimmt jedoch nicht dessen vollständige Funktionalität.

%--------------------------------------------------
\section{Randbedingungen}
\begin{itemize}
    \item Verwendung von BPMN 2.0 zur Modellierung
    \item Einsatz von Camunda als Workflow-Engine
    \item Entwicklung eines Prototyps
    \item Begrenzte Projektlaufzeit
\end{itemize}

%--------------------------------------------------
\section{Akzeptanzkriterien}
\begin{itemize}
    \item Automatischer Ablauf des Prozesses ohne manuelle Eingriffe
    \item Korrekte Verarbeitung aller Kundenreaktionen
    \item Erfolgreiche Modellierung und Ausführung in Camunda
    \item Funktionierendes Frontend zur Prozesssteuerung
\end{itemize}

%--------------------------------------------------
\section{Anhang}
Weitere Informationen.

\end{document}